\section{Identificación del fenómeno de estudio}
\label{sec:fenomeno}

\subsection{Área disciplinar y objeto tecnológico}

El estudio se inscribe en la Ingeniería Civil Informática, y de manera
específica en la intersección de tres áreas disciplinares: la inteligencia
artificial aplicada a procesos industriales, el procesamiento digital de datos
provenientes de instrumentación de planta, y el desarrollo de arquitecturas de
software orientadas a la transformación digital bajo el paradigma de la
Industria 4.0.

El objeto tecnológico central de análisis es un módulo de software analítico
para la estimación óptima de estados, concebido para operar como núcleo
matemático de consistencia física dentro de un gemelo digital. No se trata, por
tanto, de un dispositivo ni de una plataforma comercial, sino de un componente
algorítmico de la capa de software: recibe flujos de telemetría degradada,
reconstruye el estado físico latente del activo y entrega esa reconstrucción a
las capas superiores de simulación y decisión.

\subsection{Ámbito de aplicación y contexto de uso}

El ámbito de aplicación es el chancado primario en la minería extractiva del
cobre, primera etapa crítica de conminución del mineral proveniente del
yacimiento. En una estación de chancado primario giratorio de la gran minería a
cielo abierto, camiones de extracción de gran tonelaje descargan mineral de
forma discontinua sobre una tolva de recepción, desde la cual el chancador
extrae material de manera continua. El volumen retenido en esa cámara es la
variable operacional que gobierna la continuidad del proceso: un nivel demasiado
bajo produce operación en vacío e ineficiencia energética, y un nivel demasiado
alto deriva en atoros y detenciones no programadas.

La supervisión de este subsistema descansa en arquitecturas SCADA que capturan
variables operacionales mediante instrumentación industrial, típicamente
transductores ultrasónicos de nivel, sensores de potencia del motor principal y
celdas de presión. Los operadores y mantenedores de planta toman decisiones a
partir de lo que esas consolas reportan.

\subsection{Condiciones en que ocurre el fenómeno}

El entorno operativo real de una faena minera se caracteriza por condiciones
físicas extremas. La trituración de roca genera vibraciones mecánicas continuas
de gran magnitud transmitidas a toda la estructura, y el proceso libera una alta
concentración de polvo abrasivo en suspensión que interfiere directamente con la
propagación de las señales acústicas empleadas por los transductores de nivel.

Estas perturbaciones degradan sistemáticamente las señales analógicas
recolectadas e inyectan un componente de ruido estocástico de alta frecuencia
que corrompe la veracidad de la información desplegada en las consolas de
supervisión. El fenómeno es persistente y estructural: no corresponde a una
avería puntual del instrumento, sino a la condición normal de operación del
activo.

\subsection{Actores involucrados}

El fenómeno afecta a cuatro grupos de actores con intereses distintos. Los
\emph{operadores de planta} toman decisiones de alimentación y detención sobre
lecturas que no reflejan el estado real de la tolva. Los \emph{ingenieros de
procesos} calculan indicadores de rendimiento sobre series temporales
contaminadas. Los \emph{equipos de mantenimiento} pierden capacidad de
anticipación porque las señales de deterioro quedan sepultadas bajo el ruido
ambiental. Y los \emph{desarrolladores de sistemas analíticos e ingenieros
informáticos} —que son los destinatarios directos de este trabajo— construyen
modelos predictivos y tableros de decisión sobre entradas cuya calidad no está
garantizada, sin disponer de una capa intermedia que resguarde la consistencia
de los datos.

\subsection{Relación con el problema identificado}

El fenómeno tecnológico crítico radica en que las lecturas crudas de nivel
presentan desviaciones estándar que exceden los umbrales admisibles para la
automatización, introduciendo datos espurios y discontinuidades temporales.
Existe, en consecuencia, una brecha informacional severa entre el comportamiento
físico del activo y los flujos de datos que ingresan a la plataforma de
software.

La decisión analítica que orienta todo este proyecto consiste en caracterizar
dicho fenómeno como un \emph{desfase analítico de datos} y no como una falla del
hardware instrumental. Esa caracterización desplaza el foco desde la ingeniería
de instrumentación hacia la ingeniería de software, y habilita la formulación de
soluciones algorítmicas que limpien y validen la información sin intervenir la
infraestructura física existente.
